
\chapter{Introduction}
When it comes to legacy systems many problems can arise if they aren't dealt with in the correct manner,
and with the standards of technology constantly changing, and various older technologies becoming outdated or no longer supported,
modernising legacy systems is a necessary and very important step to ensuring the longevity of a system.
This dissertation aims to demonstrate the process of modernising an existing legacy database management system,
which is currently used by an organisation called BFG to keep track of fungi they have found and collected for research purposes,
and due to the use of unsupported technology, it's starting to face major problems when meeting new requirements.
The steps taken will be, first performing an analysis on the existing database and DBMS
to find any issues that need solving and to understand the new requirements of the system.
Followed by the design and implementation of the new system with more modern technologies and better practices
ensuring it meets all the new requirements.
Lastly, I will be implementing 3 variants of the same search recommendation algorithm
to demonstrate how database search could be performed whilst foraging for fungi in a place with little to no internet connection.

still to add to the introduction:
\begin{itemize}
\item list out the secific techniques covered
\end{itemize}

\section{Background}
\begin{itemize}
\item who bfg are, and what they do + a bit on FRDBI too
\item what is the study of mycology and what does it aim to achieve
\item a breif discription of bfg's database issues
\end{itemize}